% Generated draft: requires mathematical and source review.
\documentclass[11pt]{article}
\usepackage[margin=25mm]{geometry}
\usepackage{fontspec}
\setmainfont{Latin Modern Roman}
\usepackage{amsmath,amssymb,mathtools,unicode-math}
\setmathfont{Latin Modern Math}
\usepackage{enumitem,needspace,xurl,hyperref}
\setlength{\parindent}{0pt}
\setlength{\parskip}{5pt}
\setlength{\emergencystretch}{3em}
\newcommand{\N}{\mathbb{N}}
\newcommand{\Z}{\mathbb{Z}}
\newcommand{\Q}{\mathbb{Q}}
\newcommand{\R}{\mathbb{R}}
\newcommand{\C}{\mathbb{C}}
\newcommand{\F}{\mathbb{F}}
\newcommand{\A}{\mathbb{A}}
\newcommand{\PP}{\mathbb{P}}
\newcommand{\E}{\mathbb{E}}
\newcommand{\eps}{\varepsilon}
\newcommand{\dd}{\,\mathrm{d}}
\begin{document}
\section*{UnsolvedMath ID: 6200018}
\textbf{Source problem code:} catalogue entry
\subsection{Definitions and mathematical statement}
A CAT(0) group is a group $G$ admitting a geometric action, that is, a proper cocompact isometric action on a proper geodesic CAT(0) space $X$. The visual boundary $\partial_\infty X$ consists of asymptotic classes of geodesic rays and carries the cone topology.

The intended uniqueness question can be expressed as follows. Suppose $G$ does not split nontrivially over a small subgroup. For any two proper CAT(0) spaces $X,Y$ admitting geometric $G$-actions, must there be a homeomorphism
\[
\partial_\infty X\cong\partial_\infty Y?
\]
In this reading the boundary denoted $\partial_\infty G$ would be well defined up to homeomorphism. It is not presupposed to be independent of $X$.

A nontrivial splitting is ordinarily encoded by a nontrivial amalgamated product or HNN decomposition, or equivalently an action on a simplicial tree without a global fixed vertex with the relevant edge stabilizers. The source does not define ``small''. One common group-theoretic meaning is absence of a nonabelian free subgroup; slender, virtually abelian, and virtually cyclic are different possible restrictions and must not be substituted for it. Nor does the source say whether uniqueness is ordinary topological, $G$-equivariant, or a stronger metric/Tits-boundary statement. The display makes the usual visual-homeomorphism candidate explicit, pending these missing conventions.
\subsection{Short English statement}
Does forbidding splittings over small subgroups force all geometric CAT(0) models of a group to have homeomorphic visual boundaries? Specify smallness and the intended notion of boundary equivalence.
\subsection{Sources}
\begin{itemize}
\item Local UnsolvedMath registry, permanent ID 6200018, exact source code catalogue entry, ``Boundaries of Groups and Kleinian Groups---Problem 18''; proposed by Dani Wise, 2005. Registry status open, classification OPEN-TRIAGE.
\item Kapovich, \emph{Problems on Boundaries of Groups and Kleinian Groups} (2005), Problem 18, PDF page 7, \url{https://www.math.ucdavis.edu/~kapovich/EPR/problems.pdf}. Supplied normalized pdf-text extraction accessed July 29, 2026. No PDF was accessed here.
\item The supplied research summary calls the question a refinement of Croke--Kleiner phenomena and mentions related flat-torus/product results, without precise references or theorem hypotheses. It describes the general statement as open; that is a supplied assessment, not an independent status check.
\item Provenance flags current status and statement-reuse terms NEEDS\_REVIEW; the source is publicly author-hosted.
\end{itemize}
\end{document}
